\chapter{Start with the report}
\label{ch:report-first}
The report is the priority map for this edition. The reference is the unsuffixed
\path{23539195_BrauneDR_FINALREPORT.pdf} in Downloads (Prism creation date
3 November 2026, 108 PDF pages). Its 35 numbered equations, named algorithms,
literature comparisons and appendix calculations are indexed below. The current
implementation reference is the captured \path{Software/DemoCode} tree. A formula
describing the report's experiment is not automatically a statement about every
current request. Historical results remain historical; the deployed ODROID's
files, environment and model arrays were not inspected or executed.

\Q{Purpose}{Read the report methods before the helper catalogue}
Read the report equation cards here, the longer explanation at the start of
each corresponding module chapter, and its method-family comparison in
Chapter~\ref{ch:methods}. The detailed function sections then connect this theory
to current names and signatures. Small file/path helpers receive shorter notes.
Training is included because its losses, gradients and validation appear in the
report, while being clearly outside the live inference path.

\Q{Evidence}{What a reason and a source establish}
A published method citation establishes a definition or method origin.
\textbf{Derived} arithmetic follows from stated assumptions; it does not measure
those assumptions. \textbf{Implementation evidence} is the captured source.
\textbf{Reported observation} is a supplied record that was not repeated here.
\textbf{Design rationale} below explains an engineering trade-off; unless a
record is identified it is not an invented account of why a historical choice
was made. Exact heuristic coefficients have no fabricated literature optimum.

\section{Report equations one through eleven}
\ReportEq{1}{Belt travel and pitch radius}
{$\ell_{\rm rev}=N_tp=40\,\mathrm{mm/rev}$ and
$r=\ell_{\rm rev}/(2\pi)=6.366\,\mathrm{mm}$.}
{Translate motor rotation to belt motion and establish the moment arm for torque.
The current configuration uses 20 teeth, 2 mm pitch and 40 mm per revolution.
The pitch radius is not necessarily the pulley outer radius.}
{GantryConfig and its steps-per-mm properties. Geometry derived from the stated
transmission; motor/driver relations are conditional on actual hardware.}
{a4988datasheet,code2026}
\ReportEq{2}{Linear speed to motor speed}
{$n=v/\ell_{\rm rev}$, so $4000/40=100$ revolutions per minute.}
{Compare a proposed traverse rate with a speed-dependent motor curve. Millimetres
per minute divided by millimetres per revolution yields revolutions per minute.
It is a design calculation, not a measured maximum feed.}
{Mechanical sizing and \_GcodeLinearMove feed interpretation.}
{code2026,openstaxrotation}
\ReportEq{3}{Torque curve interpolation}
{$T(n)=T_1+(n-n_1)(T_2-T_1)/(n_2-n_1)$.}
{Estimate between two verified curve samples. The report's $(60,7)$ and
$(120,21)$ RPM/N\,cm points give 16.33 N\,cm at 100 RPM, but its curve citation
is literally left as SOURCE. The arithmetic is reproducible; the points and
their operating conditions remain unauthenticated.}
{Report hardware justification, not a torque lookup in the current controller.}
{finalreport2026,openstaxrotation}
\ReportEq{4}{Steady opposing load}
{$T=r\mu mg=0.0002373\,\mathrm{N\,m}=0.02373\,\mathrm{N\,cm}$ under the report inputs.}
{Estimate one steady-friction term. The assumed $\mu=0.004$ is not an observed
assembly coefficient. Acceleration, pen drag, belt forces and rotary inertia
are omitted, so this cannot establish the available dynamic margin.}
{Report mechanical sizing; the fuller load discussion is in the motion chapter.}
{openstaxrotation,finalreport2026}
\ReportEq{5}{Command lattice and rounding}
{$K=200(16)/40=80\,\mathrm{steps/mm}$ and $\Delta=1/K=0.0125$ mm.}
{Relate integer pulse targets to millimetres. Nearest rounding gives at most
$\Delta/2$ per axis and $\sqrt2\Delta/2$ endpoint distance for equal axes.
These are quantisation bounds, not loaded microstep accuracy or measured
reproduction precision.}
{GantryConfig, Calibrate and \_GcodeLinearMove.}
{a4988datasheet,code2026}
\ReportEq{6}{Driver current limit}
{$I_{\rm trip,max}=V_{\rm REF}/(8R_S)$; with $R_S=0.1\,\Omega$ and 1.68 A,
$V_{\rm REF}=1.344$ V.}
{Set a conditional electrical current limit. Eight is the driver's internal
sense scaling, not the microstep setting. Installed sense resistance, actual
reference voltage and thermal conditions must be established separately.}
{Hardware setting, not a software-assigned phase current.}
{a4988datasheet,pololu4988}
\ReportEq{7}{Minimum base resistance}
{$R_{B,\min}=(V_{\rm GPIO}-V_{BE})/I_{\rm GPIO,max}=650\,\Omega$ for the stated assumptions.}
{Limit GPIO base current. This is a minimum resistance, not a maximum. The
report's 4 mA GPIO capability lacks a resolved specific source, and 0.7 V is
an approximation; neither is certified by the arithmetic.}
{Pen transistor circuit and PEN\_MOTOR\_PIN output.}
{rohmnpn,bc337data,finalreport2026}
\ReportEq{8}{Base current and saturation check}
{$I_B=(3.3-0.7)/1000=2.6$ mA.}
{Check the selected resistor against the stated GPIO budget and transistor
drive. A forced gain of ten corresponds to 26 mA, below the listed 30 mA
motor rated current and far below its 230 mA stall current. Active-region gain
does not prove saturation.}
{Pen switching, not a pressure controller.}
{rohmnpn,bc337data,finalreport2026}
\ReportEq{9}{Switch pull down current}
{$I=V/R=3.3/10000=0.33$ mA.}
{Explain the closed-switch load and why an open logic node needs a defined
level. Leakage and noise sensitivity constrain resistance too; the result
does not prove that 10 k$\Omega$ is optimal.}
{Endstop and pen feedback wiring.}
{kernelgpio,code2026}
\ReportEq{10}{Measured axis scale}
{$K_{\rm axis}=N/L$ and $x=n/K-5$ mm in the clearance-offset frame.}
{Use counted travel between contacts and an external length reference to
establish position. The 5 mm offset reserves switch clearance; it is not a
measurement of switch repeatability.}
{Calibrate and \_GcodeLinearMove use per-axis calibrated factors.}
{code2026,bresenham1965}
\ReportEq{11}{Bresenham minor step count}
{$m_k=\left\lceil[k a_y-\lfloor a_x/2\rfloor]/a_x\right\rceil$ for the stated
major-axis accumulator convention.}
{Explain how integer error updates distribute minor-axis pulses along a line.
Initial error and the strict negative test determine ties. This is spatial
interpolation; sequential pulses and sleeps affect timing separately.}
{BresenhamMove.}
{bresenham1965,code2026}

\section{Report equations twelve through twenty four}
\ReportEq{12}{Integral rectangle sum}
{$Q=S(y_1,x_1)-S(y_0,x_1)-S(y_1,x_0)+S(y_0,x_0)$.}
{Obtain a window total from four prefix-sum entries using inclusion and
exclusion. The full table still costs image-sized work; constant window work
does not mean constant full-image processing time.}
{Integral, BoxSum and BoxMean.}
{crow1984}
\ReportEq{13}{Background ratio and ink mask}
{$R=G/\max(B_r,1)$ and $M=\mathbf1[R<T]$.}
{Reduce multiplicative illumination changes before thresholding. The floor
prevents zero division; glare, clipping and ink-filled windows can still
invalidate the paper-background approximation.}
{InkMask and SegmentLines crop normalisation.}
{code2026,otsu1979}
\ReportEq{14}{Otsu threshold criterion}
{$T=\arg\max_t\omega_0(t)\omega_1(t)[\mu_0(t)-\mu_1(t)]^2$.}
{Separate the histogram into two maximally distinct classes. Histogram
optimality is not semantic optimality for true ink, and empty classes require
guards. Bins, range, tie handling and comparison polarity matter.}
{Otsu, InkMask and BinarizeLine.}
{otsu1979,opencvthreshold}
\ReportEq{15}{Deskew concentration}
{$\theta^*=\arg\max_\theta\sum_y p_\theta(y)^2$ with
$p_\theta(y)=\sum_x M_\theta(y,x)$.}
{Prefer rotations concentrating horizontal text into fewer rows. The selected
angle grid is a project search budget, not a general guarantee on curved or
perspective-distorted writing.}
{SkewScore and Deskew; do not confuse with slant estimation.}
{postl1986}
\ReportEq{16}{Local seam cost}
{$C(y,x)=4M_c(y,x)+0.012[(y-y_m)/h_g]^2$.}
{Avoid closed-mask ink while weakly attracting the boundary toward the middle
of its corridor. The dimensionless weights encode project priorities; the
literature does not establish these exact coefficients.}
{SeamCuts.}
{arvanitopoulos2014,code2026}
\ReportEq{17}{Seam recurrence}
{$D(y,x)=C(y,x)+\min_d[D(y+d,x-1)+0.02|d|]$.}
{Compare complete legal prefixes rather than greedily selecting one local
pixel. Backtracking recovers the route; the optimum is limited to the chosen
corridor and transitions. SeamCuts explicitly calls step 7 despite the
LeastCostSeam default 3.}
{LeastCostSeam and SeamCuts.}
{bellman1954,code2026}
\ReportEq{18}{Fixed canvas scale}
{$a=\min(64/h,640/w)$, with resized dimensions near $(ah,aw)$.}
{Preserve aspect ratio while fitting the checkpoint input. Long lines can
become short vertically when width determines scale. Padding is vertically
centred and left aligned.}
{FitFixedCanvas, ResizeLineImageFixed and FrameXToPixel.}
{wolberg1990,code2026}
\ReportEq{19}{Convolution response}
{$Z_{o,y,x}=b_o+\sum_{c,i,j}W_{o,c,i,j}X_{c,y+i,x+j}$ on the appropriately padded array.}
{Extract local learned stroke patterns shared across positions. This is
cross-correlation in common framework terminology. A feature channel is not
automatically a labelled character detector.}
{Conv2d and the train-mode convolution pair.}
{shi2017,torchconv}
\ReportEq{20}{Batch normalisation}
{$\hat z=(z-\mu)/\sqrt{v+\epsilon}$ and $y=\gamma\hat z+\beta$.}
{Control activation scale while allowing a learned affine response. Training
uses current batch statistics; inference uses saved running statistics.
Changing that mode changes the operation.}
{BatchNorm2d and BatchNorm2dForwardTrain/Backward.}
{ioffe2015,torchbn}
\ReportEq{21}{LSTM cell and hidden state}
{$c_t=f_t\odot c_{t-1}+i_t\odot g_t$ and $h_t=o_t\odot\tanh c_t$.}
{Carry gated memory along feature columns and expose a bounded hidden output.
The gates and candidate are computed from current features and previous
hidden state; weights are shared, not passed as the state.}
{LstmForward and LstmBackward.}
{hochreiter1997,torchlstm}
\ReportEq{22}{CTC sequence loss}
{$L=-\log\sum_{\pi:B(\pi)=\ell}\prod_t p_t(\pi_t)$.}
{Learn complete transcriptions without character boxes by summing all valid
alignment paths. Collapse merges adjacent repeats before removing blank;
literal spaces remain ordinary labels.}
{CtcLoss; DecodeCtc solves a cheaper best-path inference problem.}
{graves2006,ctcdistill}
\ReportEq{23}{CTC occupancy and derivative}
{$q_{t,k}=\sum_{s:\ell'_s=k}\exp(\alpha_{t,s}+\beta_{t,s}-\log Z)$ and
$\partial L/\partial z_{t,k}=p_t(k)-q_{t,k}$.}
{Propagate the sequence loss through logits. Alpha includes the current
emission and beta excludes it in this convention, preventing double counting.
State indices differ from class indices.}
{CtcLoss; batch averaging and library reduction must also match.}
{graves2006,torchctc}
\ReportEq{24}{Controlled stroke width}
{$I_s=255[1-\delta_k(\mathrm{thin}(M))]$,
$w=\mathrm{clip}(\mathrm{round}(\rho h_x),2,9)$,
$k=2\lfloor w/2\rfloor+1$.}
{Reduce pen-width cues before shape-based identification. The effective odd
kernel quantises the requested relative width; thinning can change geometry.
Current PenWidthXh is 0.17. Historical sweep selection is not an independent
qualification of this source snapshot.}
{StrokeNormalize and ClassifyImage.}
{zhangsuen1984,finalreport2026,code2026}

\section{Report equations twenty five through thirty five}
\ReportEq{25}{Forced transcript alignment}
{$V(t,s)=\log p_t(\ell'_s)+\max_{r\in\mathcal P(s)}V(t-1,r)$.}
{Recover a high-scoring legal path consistent with the known label before
extracting characters. Replacing path summation by maximisation gives a hard
alignment, not proof of correct geometric boundaries.}
{CtcForcedAlign and ExtractLineGlyphs.}
{graves2006,viterbi1967}
\ReportEq{26}{Character cut cost}
{$C_x=p_x/q_{0.60}(p_j:p_j>0)+1.15\,\mathrm{clip}[(b-y_{\rm top}(x))/h_x,0,3]$.}
{Prefer thin low connectors over stems near the recogniser's suggested cut.
Blank columns have zero cost. Ordered boundaries and minimum cell width
prevent neighbouring cuts collapsing. All exact weights/percentiles are
project heuristics, not derived from CTC.}
{ExtractLineGlyphs.}
{code2026}
\ReportEq{27}{Stored glyph coordinates}
{$v=(b-y)/h_x$, $u_0=(x-x_L)/h_x-v\tan\phi$, and $u=u_0-\delta$.}
{Flip image-down to glyph-up, remove source scale and slant, and store trim
as lead spacing. Shear is not rotation. The inverse placement must restore
the same conventions.}
{ExtractLineGlyphs and SynthesizeText.}
{wolberg1990,vinciarelli2001,code2026}
\ReportEq{28}{Measured connectedness}
{$c=N_{\rm cross}/\max(1,N_{\rm adj})$.}
{Summarise eligible in-word boundaries crossed by a traced centreline. It is
not measured original pen-lift probability: touching and bad cuts can inflate
it. Current synthesis suppresses joins below 0.45 and chooses cursive fallback
style at 0.5.}
{BuildAuthorProfile, \_CoreGlyph and SynthesizeText.}
{code2026,bulacu2007}
\ReportEq{29}{New glyph placement}
{$x=p_x+H(l+u+sv)$ and $y=p_y+H(v+d)$.}
{Restore writer shear, physical size, lead and baseline drift at the new
text location. H and origins have millimetre units; stored glyph coordinates
are dimensionless. Additional jitter/collision handling changes placement.}
{SynthesizeText.}
{code2026,wolberg1990}
\ReportEq{30}{Automatic anchoring strength}
{$\lambda=\mathrm{clip}[(d-0.24)/(0.32-0.24),0,1]$.}
{Increase fallback preference when median stored prior distance is poor.
The endpoints are project settings, not universal perceptual thresholds.
A stored explicit legibilityLambda can override automatic calculation.}
{AutoLegibility, \_SwapProb and SynthesizeJointBestOf.}
{code2026}
\ReportEq{31}{Candidate objective}
{$A=\max(0,1-d_{\rm edit}/|t|)$ and $J=2AP/(A+P)$ for $A+P>0$.}
{Prefer candidates with both text recovery and requested-writer probability.
The harmonic mean penalises a weak component, but P is a model posterior
and A is an edit score, so J is not a joint likelihood or independent physical
style judgement.}
{SynthesizeJointBestOf and CharAcc.}
{code2026,levenshtein1966}
\ReportEq{32}{Machine placement}
{$x_m=o_x+m_x+a(x-x_0)$,
$y_m=o_y-m_y-d_y-a(y_1-y)+\Delta y$.}
{Place the trajectory within margins and fit to bounds. The current
ToMachine retains scaled synthesis pitch and removes measured row-centre
wander; it has no ignoreRuling choice. Both frames here are y-up.}
{ToMachine and \_Clamp.}
{code2026,wolberg1990}
\ReportEq{33}{Nominal timing budget}
{$t_{\rm nom}=L_d/v_d+L_u/v_u+N_{\rm toggle}\tau$.}
{Explain drawing, travel and dwell contributions. Current WriteGcode uses
feeds in mm/min, multiplies length/feed by 60, and approximates dwell using
the mean of up/down dwell. Initial/final travel, actual toggles and software
overhead prevent treating it as a rigorous measured lower bound.}
{WriteGcode and timing.Run; keep estimates distinct from completed-job logs.}
{code2026,pythontime}
\ReportEq{34}{Pooled character recovery}
{$A_c=100[1-\sum_i d_i/\sum_i|t_i|]$.}
{Weight characters rather than equally weighting lines. This differs from
the clamped per-candidate CharAcc at extreme insertion counts. Case, spaces,
punctuation and empty-reference handling must remain explicit.}
{Report evaluation and server scoring.}
{levenshtein1966,finalreport2026}
\ReportEq{35}{Geometric agreement and natural baseline}
{$F=2PR/(P+R)$ and $Q_k=100F_{{\rm robot},k}/F_{{\rm genuine},k}$.}
{Compare thresholded centreline precision/recall and then contextualise robot
agreement using genuine writing variation. F is not character accuracy or IoU;
Q is project-defined and can exceed 100 percent. Mean writer ratios differ
from a ratio of pooled means.}
{Historical geometric evaluation, not a loaded classifier confidence.}
{tanksmetric,finalreport2026}

\section{Other named report methods that must not disappear}
\Q{Method}{Normalisation and affine terminology}
Min--max scaling uses $(x-x_{\min})/(x_{\max}-x_{\min})$ to map a range;
z-score scaling uses $(x-\mu)/\sigma$ to centre and scale a distribution.
Neither is identical to illumination division, stroke-width normalisation or
geometric x-height normalisation. Each removes a different nuisance variable.
The denominator needs a defined constant-image convention. These are direct
definitions, not evidence that all are active inference stages.

\Q{Alternatives}{LeNet five and residual networks}
The report mentions LeNet-5 as a historical convolution/pooling example and
residual networks as a deeper architecture with shortcut paths. A residual
block computes $y=x+F(x)$ (or a projection when dimensions differ); the extra
route can help gradient propagation. Neither is the current twelve-convolution,
two-BiLSTM reader. Changing the backbone would require compatible weights and
new evaluation; citing a benchmark architecture does not establish better
accuracy or speed for this photographed dataset~\cite{lenet1998,resnet2016}.
Its historical radial-basis output compares feature position with learned
class representatives through a distance-based response; it is not the current
affine writer head or framewise text softmax. ANN is the broad neural-network
family, CNN specifies local shared spatial weights, and RNN specifies a state
carried along a sequence. These names are not separate active modules merely
because they appear in the literature chapter.

\Q{Alternatives}{Recurrent stroke generation GANwriting and One DM}
Graves's recurrent synthesis predicts successive online movements; the project
has offline images with no observed stroke order. GANwriting conditions raster
generation on content and writer style; its image still needs a pen-path
conversion for this robot. One-DM iteratively denoises an image and can use one
style reference at generation, but it was trained on a larger corpus. A small
reference requirement is not a small training-data requirement. The current
glyph library avoids training a separate generator and exposes controllable
geometry, while finite variants approximate contextual cursive poorly. These
are report-grounded trade-offs, not new head-to-head measurements
~\cite{gravesgenerate,ganwriting,onedm,haines2016}.

\Q{Method}{Writer texture allographs and ink width}
Textural edge-direction distributions summarise local stroke orientations;
allographs are alternative realisations of a character. The explicit glyph
pool stores allographic variation, while the CNN learns discriminative maps.
Ink width can carry writer evidence but changes when the robot uses one pen.
That motivates shape normalisation, without proving all ink information is
undesirable or every normalised prediction is more accurate
~\cite{bulacu2007,brink2012}.

\Q{Alternatives}{Mechanisms lifts and control}
A Cartesian gantry maps orthogonal axes directly; serial arms require coupled
kinematics and can vary stiffness over the workspace; parallel mechanisms
trade workspace and coupling differently. The report's one-way geared cam
replaces a position-commanded servo or third stepper with a switched motor and
binary raised-state feedback. A4988 current chopping controls winding current;
it is not the same as the controller's STEP sequence. Open-loop interior motion
and switch-referenced homing are not continuous position feedback. These
mechanical trade-offs are discussed here because they determine the meaning of
the software calculations, not because another mechanism was built or tested.

\Q{Evidence}{Appendix calculations plots and recorded checks}
The two-crop page vote, finite-difference derivative, patch-storage calculation,
12 ms switch confirmation, stroke-width sweep, confusion matrices, seam/capture
counts and training epoch logs are all retained in the verification chapter.
For a confusion matrix, row normalisation gives prediction fractions conditional
on actual writer; it is not class precision. Wilson limits describe binomial
decision counts subject to dependence caveats. A recognition score used to
select a candidate is not an independent physical test of that candidate.
Every retained historical number stays attached to its population and source.

\section{Report algorithm and figure reading map}
\Q{Purpose}{The named procedures and their current homes}
The pen-toggle procedure is \texttt{SetPen}; debounced switch reads are
\texttt{TriggeredEndstop}; dynamic axis calibration uses
\texttt{\_CalibrateAxis} and \texttt{Calibrate}; coordinated integer movement
is \texttt{BresenhamMove}. The report's line-ownership refinement maps to
\texttt{AssignOwnership} and its correction helpers. Labelled handwriting to
a writer profile maps to \texttt{ExtractLineGlyphs},
\texttt{ExtractAuthorRaw} and \texttt{BuildAuthorProfile}.
Joint search/accepted repair is \texttt{SynthesizeJointBestOf}. The extended
training algorithm maps to the train-mode forward/cache/backward functions,
loss routines, optimiser step, validation and checkpoint logic. These names
are linked through the function index; unrelated helpers do not outrank them.

\Q{Method}{What each report plot can and cannot show}
The integral/ratio example demonstrates arithmetic for one pixel, not all
lighting conditions. A seam/cost overlay shows the route preferred by the
chosen objective; a labelled crop establishes whether that route preserves
the intended ink. A convolution feature plot shows activation responses,
not fixed semantic labels for channels, and independent colour scales cannot
be compared as absolute strengths. A CTC plot places symbols against feature
frames, not physical pen time. Forced alignment is conditional on the supplied
transcript. A Bresenham staircase represents integer commanded positions,
not an observed vibration trace. Stroke-normalisation panels show mask,
skeleton and restored width rather than guaranteed shape preservation.

Variant clouds, medoids and prototypes have different provenance: a medoid
is a real selected sample, whereas a prototype is constructed. Candidate-score
panels show optimisation against the judging networks. Confusion matrices
require the actual per-class counts. Word overlays depend on registration,
normalisation and distance tolerance; they are not automatic proof of
character correctness. Timing bars require host, stage boundaries and
completed-job status. Training curves require the selected checkpoint and
holdout definition. The verification chapter explains the exact reported
calculations and their qualifications rather than rerunning them.

\Q{Evidence}{System decisions mentioned outside the algorithms}
The camera, frame/clamps, Cartesian mechanism, geared lift, endstops, transistor
and power/driver connections influence assumptions in the software and are
covered in the motion/hardware explanation. Ergonomics, cost, maintainability,
privacy and environmental discussions are report design context rather than
additional image-processing functions. They do not create undocumented
sensor feedback, trained classes or software safety guarantees.
